\documentclass[12pt, a4paper]{article}

%% ── Packages ──────────────────────────────────────────────────────────────────
\usepackage[utf8]{inputenc}
\usepackage[T1]{fontenc}
\usepackage{lmodern}
\usepackage{microtype}
\usepackage[margin=2.5cm]{geometry}
\usepackage{setspace}
\usepackage{graphicx}
\usepackage{amsmath, amssymb}
\usepackage{booktabs}
\usepackage{array}
\usepackage{hyperref}
\usepackage{url}
\usepackage{natbib}
\usepackage{caption}
\usepackage{subcaption}
\usepackage{float}
\usepackage{enumitem}
\usepackage{xcolor}
\usepackage{lineno}

\hypersetup{
    colorlinks = true,
    linkcolor  = blue!60!black,
    citecolor  = green!50!black,
    urlcolor   = blue!70!black,
}

\doublespacing

%% ── Title block ───────────────────────────────────────────────────────────────
\title{\textbf{LeafSegmenter: An Interactive, GPU-Optional Desktop Application
for Plant Image Segmentation and Phenotyping Using the Segment Anything
Model~2}}

\author{
  Nirwan Tandukar\textsuperscript{1} \\[4pt]
  \textsuperscript{1}[Your Institution / Department]\\
  \texttt{[your.email@institution.edu]}
}

\date{\today}

%%──────────────────────────────────────────────────────────────────────────────
\begin{document}

\maketitle
\linenumbers

%% ── Abstract ──────────────────────────────────────────────────────────────────
\begin{abstract}
Accurate segmentation of plant structures in field and greenhouse imagery is essential for high-throughput automated phenotyping. Existing tools often require substantial computing resources, command-line expertise, or support only one segmentation approach. We present \textbf{LeafSegmenter}, an open-source, cross-platform desktop application that combines the Segment Anything Model~2 (SAM2) with interactive mask-editing, image-enhancement, and model-training tools usable by non-experts on standard laptops. The application supports GPU-accelerated inference (CUDA and Apple MPS) and CPU-only operation, avoiding the need for high-performance computing. Key features include an interactive colour-based mask filter, geometric and species-specific leaf-completion templates, a leaf-unfolding module based on mirror reflection across the detected fold line, an undo history, and a training pipeline that distils curated SAM2 masks into a lightweight U-Net for fast, automatic segmentation of new images. Comprehensive phenotyping metrics (area, length, width, shape descriptors, colour statistics, and vegetation indices) are exported as structured CSV files. LeafSegmenter lowers the barrier to plant phenotyping research while remaining extensible to larger models and new species.
\end{abstract}

\textbf{Keywords:} plant phenotyping, image segmentation, Segment Anything
Model, vision transformer, interactive tools, leaf morphology, field imaging.

\newpage

%%──────────────────────────────────────────────────────────────────────────────

\section{Introduction}
\label{sec:introduction}

Plant phenotyping is the quantitative measurement of plant morphology, architecture, color, and growth-related traits. It is central to modern crop improvement, stress physiology, and functional plant biology \citep{walter2015plant,furbank2011phenomics}. As genotyping has become faster and cheaper, precise, high-throughput trait measurement has become the main bottleneck in linking genotype, environment, and phenotype
\citep{pound2017deep,singh2016machine}. Digital imaging partly addresses this by enabling non-destructive measurements in controlled environments, greenhouses, field plots, and UAV platforms, but the bottleneck has shifted from image acquisition to image interpretation. Accurate segmentation of plant organs from backgrounds, soil, trays, shadows, neighboring plants, and overlapping tissues remains essential for reliable trait extraction.

A segmentation model can be difficult to train because plants vary widely across species, genotypes, developmental stages, imaging platforms, and environments. Leaves can overlap, fold, curl, or occlude each other; field images often have uneven illumination and complex backgrounds; and color differences between target tissues and other objects, such as grass and weeds, can be grass and subtle. Classical methods based on color thresholding, vegetation indices, and morphological operations can work under controlled conditions but often need extensive parameter tuning and fail when lighting, background, or morphology changes \citep{woebbecke1995color,gitelson2002novel}. Deep learning methods, such as U-Net-style semantic segmentation and Mask R-CNN-style instance segmentation, have greatly improved plant segmentation \citep{ronneberger2015unet,he2017mask,pound2017deep}, but supervised models still require large annotated datasets, task-specific training, and
computational expertise, limiting their use in small laboratories, early-stage projects, and scenarios with scarce annotated data.

Visual foundation models create new opportunities for plant phenotyping. Trained on large, diverse datasets, they can often be adapted to new tasks with little or no task-specific training \citep{bommasani2021opportunities}. The Segment Anything Model (SAM) introduced promptable segmentation, where users guide the model with points, boxes, or masks to obtain object masks in a zero-shot setting \citep{kirillov2023segment}. SAM~2 extends this to images and videos using a promptable segmentation architecture, hierarchical image encoder, and streaming memory, enabling improved interactive segmentation and refinement \citep{ravi2025sam2}. These models are useful for plant phenotyping because many studies need accurate masks but lack resources to annotate and train a new model for each species, organ, or imaging condition.

Applications of SAM to plant phenotyping has shown both promise and limitations. \citet{zhang2024adapting} combined SAM with image–text prompting for zero-shot plant recognition and phenotypic measurement without manually labeled training data. \citet{williams2024leafonlysam} built a zero-shot pipeline for potato leaf segmentation using SAM plus plant-specific post-processing. Their results show that SAM-based methods are useful when annotated data are scarce, but automatic outputs often include background objects, duplicate masks, whole-plant masks, or non-leaf structures that must be filtered before measurements are reliable. Foundation models thus reduce the annotation burden, but practical plant phenotyping still needs human-in-the-loop correction, mask filtering, quality control, and phenotype-specific measurement tools. Transformer-based vision models capture long-range spatial relationships via self-attention and now underpin modern segmentation systems such as SAM and SAM~2 \citep{dosovitskiy2021image,ravi2025sam2}. Yet they are often computationally expensive and rely on large-scale pretraining or carefully tuned training, whereas convolutional models can remain competitive when data or compute are limited \citep{chen2021visformer}. For routine phenotyping, this creates demand for workflows that pair the accuracy and flexibility of foundation models with lightweight models that can be trained from curated masks and quickly deployed on new images. Knowledge distillation and teacher–student training address this need: high-quality masks generated or corrected by a large model can train a smaller model for faster, task-specific inference \citep{hinton2015distilling,zhang2023mobilesam}.

Despite these advances, accessible software for plant-focused, interactive, foundation-model segmentation remains scarce. General annotation platforms such as LabelMe offer manual labelling but lack plant-specific enhancement, mask filtering, phenotypic measurement, and model training workflows \citep{russell2008labelme}. PlantCV is a powerful, widely used image-analysis library for plant phenotyping, but is mainly script-based and thus requires programming expertise for many tasks \citep{gehan2017plantcv}. Meanwhile, many SAM-based plant studies appear as pipelines or experimental frameworks rather than integrated desktop applications that let non-expert users load images, segment plants, edit masks, export traits, and train a lightweight model within a single environment.

We present \textbf{PlantSegmenter}, an interactive, python based GPU-linked desktop application for plant image segmentation and phenotyping with SAM~2. PlantSegmenter is not meant to make segmentation easily accessible in everyday plant biology, especially for users with limited programming skills, few annotations, or no high-performance computing. It combines SAM~2-based segmentation with plant-specific image enhancement, color-based mask filtering, interactive mask editing, several phenotype measurements, geometric leaf completion, folded-leaf reconstruction, and a training workflow that turns curated SAM~2 masks into a lightweight U-Net model. The software is available for download from the GitHub repository at: https://github.com/nirwan1265/Image-Segmentation.



%%──────────────────────────────────────────────────────────────────────────────
\section{Materials and Methods}
\label{sec:methods}

\subsection{Application Architecture}

PlantSegmenter is implemented in Python~3.11 with a \texttt{tkinter}-based graphical interface. The application is organized as a modular package (\texttt{app}) with seven functional modules:

\begin{itemize}
    \item \texttt{gui\_app.py} — main application class, layout, event handling
    \item \texttt{image\_processing.py} — pure image enhancement functions
    \item \texttt{mask\_utils.py} — mask mathematics and shape completion geometry
    \item \texttt{phenotyping.py} — phenotype measurement and CSV export
    \item \texttt{sam2\_utils.py} — SAM~2 model loading and Hydra configuration
    \item \texttt{color\_filter.py} — interactive colour-based mask filtering
    \item \texttt{tab\_leaf\_completion.py}, \texttt{tab\_leaf\_unfolding.py},
          \texttt{tab\_train.py} — training panel tab modules
\end{itemize}

The modular design separates computationally intensive functions (image enhancement, mask geometry, phenotype measurement) from the GUI, allowing independent use for batch processing and unit testing.

\subsection{Image Loading and Preprocessing}

The application accepts TIFF, PNG, JPEG, and BMP images. A built-in TIF converter converts raw field-acquisition formats to standard 8-bit PNG for segmentation. Images are loaded as RGB \texttt{uint8} NumPy arrays via Pillow and OpenCV python packages.

A rotation module uses a dial widget with preset angles (45°, 90°, 180°, 270°, $-$45°, $-$90°) and arbitrary rotation to orient images acquired at non-standard angles before segmentation.

\subsection{Image Enhancement}

A dedicated enhancement panel offers a configurable, sequential processing pipeline for plant imagery, including these operations:

\subsubsection{Green-Channel CLAHE Enhancement}

The main plant-specific enhancement applies Contrast Limited Adaptive Histogram Equalisation (CLAHE) \citep{zuiderveld1994clahe} only to the value channel (V) of pixels classified as green by HSV thresholding. This boosts contrast within leaf tissue while leaving the background unchanged, improving the signal-to-background ratio for SAM~2. A bilateral filter \citep{tomasi1998bilateral} then reduces noise while preserving leaf edges, followed by Sobel edge blending \citep{sobel1968isotropic} to highlight boundaries. An unsharp masking step \citep{dahl2008contrast} concludes the standard plant pipeline.

\subsubsection{Vegetation Indices}

The application computes five common vegetation indices from the RGB channels, usable as visual overlays or input features for downstream analysis:

\begin{itemize}
    \item \textbf{ExG} (Excess Green): $\mathrm{ExG} = 2G - R - B$
          \citep{woebbecke1995color}
    \item \textbf{ExGR}: $\mathrm{ExGR} = \mathrm{ExG} - 1.4R - G$
    \item \textbf{GRVI} (Green-Red Vegetation Index):
          $\mathrm{GRVI} = (G - R)/(G + R)$ \citep{tucker1979red}
    \item \textbf{VARI} (Visible Atmospherically Resistant Index):
          $\mathrm{VARI} = (G - R)/(G + R - B)$ \citep{gitelson2002novel}
    \item \textbf{GLI} (Green Leaf Index):
          $\mathrm{GLI} = (2G - R - B)/(2G + R + B)$ \citep{louhaichi2001spatially}
\end{itemize}

\subsubsection{Advanced Enhancement Options}

Additional enhancement methods are available for challenging imaging
conditions:

\begin{itemize}
    \item \textbf{Retinex} (single- and multi-scale) \citep{jobson1997multiscale}
          for illumination normalisation under variable field lighting.
    \item \textbf{Morphological top-hat transform} \citep{serra1983image} for
          uneven background illumination correction.
    \item \textbf{Non-local means denoising} \citep{buades2005non} for
          high-noise images.
    \item \textbf{Guided filtering} \citep{he2013guided} for edge-preserving
          smoothing using the green channel as guide.
    \item \textbf{LAB colour-space enhancement}, \textbf{grey-world and
          max-white white balance}, \textbf{Difference of Gaussians},
          \textbf{local contrast normalisation}, \textbf{adaptive gamma
          correction}, and \textbf{shadow/highlight recovery}.
    \item \textbf{Background whitening}: HSV-based detection of
          bright, desaturated background pixels (e.g.\ white imaging trays)
          with morphological cleanup, replacing them with white to simplify
          foreground segmentation.
\end{itemize}

\subsection{SAM2-Based Segmentation}

PlantSegmenter integrates SAM~2 \citep{ravi2024sam2} via the official Python API. The user selects a SAM~2 checkpoint (e.g., Hiera-Small, Hiera-Large, or a bundled \texttt{.pt} file) and the compute device (CPU, CUDA, or Apple MPS). After image enhancement, segmentation runs on the current image using \texttt{SAM2AutomaticMaskGenerator}, with all hyperparameters configurable in the GUI:

\begin{itemize}
    \item \texttt{points\_per\_side}: grid density of prompt points.
    \item \texttt{pred\_iou\_thresh}: predicted IoU quality threshold.
    \item \texttt{stability\_score\_thresh}: mask stability rejection threshold.
    \item \texttt{crop\_n\_layers}: number of multi-scale crop passes.
    \item \texttt{crop\_overlap\_ratio}: overlap between adjacent crops.
    \item \texttt{box\_nms\_thresh}: duplicate suppression threshold.
    \item \texttt{min\_mask\_region\_area}: minimum mask area in pixels.
    \item \texttt{use\_m2m}: mask-to-mask refinement flag.
\end{itemize}

Hover tooltips in the GUI also explain each parameter. After generation, masks are deduplicated using Intersection-over-Union (IoU $> 0.80$) and split into single-connected-component instances. Batch segmentation processes an entire folder, caching results in memory for later review and export.

\subsection{Mask Editing Tools}

A dedicated Masks panel provides interactive editing operations:

\begin{itemize}
    \item \textbf{Delete}: remove selected masks from the set.
    \item \textbf{Combine}: union of two or more selected masks into one.
    \item \textbf{Refine}: re-run SAM~2 within the bounding box of a selected
          mask using denser prompting for improved boundary accuracy.
    \item \textbf{Edit boundary}: an interactive contour editor allowing
          point-by-point drag adjustment of mask boundaries, with add-point
          (Ctrl+click) and delete-point (Shift+click) operations and optional
          Gaussian smoothing.
    \item \textbf{Split}: decompose multi-blob masks into individual
          connected-component masks.
\end{itemize}

An undo history system (up to 50 states) captures the full mask list before
each destructive operation, enabling recovery via Ctrl+Z (or Cmd+Z on macOS).

\subsection{Color-Based Mask Filtering}

A color filter panel enables interactive mask selection from pixel color statistics. A 2D Hue × Saturation picker renders the HSV space as a gradient, where users drag rectangles to define color ranges combined with OR logic. An eyedropper samples pixels from the preview canvas, places a crosshair on the picker, and auto-creates a selection rectangle with suitable tolerance. A global brightness (Value) range control applies to all selections. The filter uses the mean HSV color of each mask’s pixels, allowing fast removal of soil, background, or specular highlight masks without manual selection.

\subsection{Leaf Shape Completion}
\label{sec:completion}

Many plant images contain partially occluded or truncated leaves. To enable complete leaf reconstruction for morphometric analysis, PlantSegmenter includes a leaf completion module with two complementary shape galleries:

\subsubsection{Geometric Shape Gallery}

Eighteen botanical leaf forms—Orbicular, Elliptical, Ovate, Obovate, Cordate, Reniform, Lanceolate, Oblanceolate, Oblong, Linear, Ensiform, Cuneate, Spathulate, Deltoid, Rhomboid, Sagittate, Hastate, and a general Ellipse—are provided as silhouette thumbnails defined by parametric polygons.

\subsubsection{Crop Species Gallery}

Eighteen crop-specific silhouettes are provided, traced from documented leaf morphology: Maize, Wheat, Sorghum, Rice, and Barley (narrow strap-leaf monocots); Arabidopsis (rosette); Tomato and Potato (pinnate compound); Soybean (trifoliate); Sunflower, Tobacco, Cassava (palmate); and Cotton, Grape, Common Bean, Cucumber, Pepper, and Lettuce.

\subsubsection{Alignment Algorithm}

Given a selected mask and shape template, the completion algorithm performs
three sequential alignment steps:

\begin{enumerate}
    \item \textbf{Centroid alignment}: the shape template is translated so that
          its centre coincides with the pixel centre of mass (centroid) of the
          mask.
    \item \textbf{PCA orientation}: Principal Component Analysis is applied to
          the mask pixel coordinates.  The major eigenvector defines the leaf's
          primary axis; the template is rotated to align its major axis with
          this direction.
    \item \textbf{Scale fitting}: the template axes are scaled to match the
          PCA-derived major and minor axis lengths of the mask.
\end{enumerate}

The aligned template is rasterised at full resolution and merged with the original mask (preserving all pixels). Fine interactive adjustments to scale, position, and rotation are available before committing the final mask.

\subsection{Leaf Unfolding Reconstruction}
\label{sec:unfolding}

Folded leaves are common in greenhouse and field imagery, especially for maize, sorghum, and wheat. A Leaf Unfolding module reconstructs the full leaf by mirroring the folded part across an automatically detected fold line.

The fold line is found as the boundary between two selected masks representing the folded halves. Both masks are morphologically dilated, and their intersection defines a junction zone. Singular Value Decomposition (SVD) on the junction pixel coordinates yields the fold line direction as the first right singular vector. Each pixel $\mathbf{p}$ of the folded mask is then reflected as:

\begin{equation}
    \mathbf{p}' = \mathbf{p} - 2\left[(\mathbf{p} - \boldsymbol{\mu}) \cdot
    \hat{\mathbf{n}}\right]\hat{\mathbf{n}}
    \label{eq:reflection}
\end{equation}

where $\boldsymbol{\mu}$ is a point on the fold line and $\hat{\mathbf{n}}$ is its unit normal. The reflected mask is morphologically closed to remove rounding artefacts, and the corresponding image pixels are mapped to their reflected locations. Before applying, the user can interactively adjust the fold line angle (in $\pm$5° steps) and shift the mirrored piece.

\subsection{Phenotype Measurement and Export}

For each mask (or set of selected masks), the following measurements are
computed and exported as CSV files:

\subsubsection{Morphometric Descriptors}
\begin{itemize}
    \item \textbf{Area} (px$^2$)
    \item \textbf{Length and width}: PCA major/minor axis lengths; deskewed
          row-span length and 95th-percentile row width.
    \item \textbf{Perimeter}, \textbf{convex hull area and perimeter}
    \item \textbf{Solidity}: mask area / convex hull area
    \item \textbf{Circularity}: $4\pi A / P^2$
    \item \textbf{Extent}: mask area / bounding box area
    \item \textbf{Equivalent diameter}: diameter of a circle of equal area
    \item \textbf{Connected components}: number of distinct blobs
\end{itemize}

\subsubsection{Colour Descriptors}
\begin{itemize}
    \item Per-channel (R, G, B, H, S, V) mean, median, sum, and standard
          deviation within the mask
    \item HSV channel variances
    \item Vegetation index means: ExG, ExR, ExGR, GLI, green pixel fraction
\end{itemize}

Both \textit{individual} (per-segment) and \textit{joint} (aggregated across
a selection) CSV exports are supported.  Masks may also be saved as full-size
binary PNG files, bounding-box-cropped binary masks, and RGBA transparent-background
crops.

\subsection{Custom Model Training}
\label{sec:training}

To enable rapid, class-specific segmentation without SAM~2 at inference time,
PlantSegmenter includes an integrated training pipeline.

\subsubsection{Data Collection}

Users gather training examples during normal SAM~2-assisted segmentation: after segmenting an image, selected masks are added to a structured dataset (images in \texttt{images/}, binary masks in \texttt{masks/}) with one button press. Images can also be labeled as negative examples (no target present).

\subsubsection{Model Architecture}

The default architecture is a U-Net \citep{ronneberger2015unet} with a ResNet-18 \citep{he2016deep} encoder pre-trained on ImageNet \citep{deng2009imagenet}, implemented using \texttt{segmentation-models-pytorch} \citep{yakubovskiy2019smp}. A lightweight variant (\texttt{unet\_small}) omits the pre-trained encoder for resource-constrained settings. The binary decoder outputs a single-channel mask.

\subsubsection{Training Procedure}

Training uses the Adam optimiser \citep{kingma2014adam} with configurable learning rate (default $10^{-5}$), batch size, image size, device, and gradient steps. A weighted binary cross-entropy loss addresses foreground–background class imbalance. Training can be resumed from a checkpoint, and all hyperparameters are configurable via the GUI without code changes.

\subsubsection{Inference}

The trained model accepts arbitrary-resolution images, runs the same enhancement pipeline as SAM~2 segmentation, and generates instance masks via connected-component analysis on a thresholded probability map. A sliding-window inference mode processes large images using overlapping tiles and merges predictions, while an optional colour-guided mode restricts tiles to plant-coloured regions to substantially reduce inference time on large field images.

%%──────────────────────────────────────────────────────────────────────────────
\section*{Results}

[Your acknowledgements here.]

%%──────────────────────────────────────────────────────────────────────────────
\bibliographystyle{abbrvnat}
\bibliography{references}

\end{document}
